\section{Introduction}
\label{sec:introduction}

Laser-Induced Breakdown Spectroscopy (LIBS) is an atomic emission technique in which a focused laser pulse ablates a small volume of sample material, generating a transient plasma whose optical emission encodes the elemental composition of the target.
Calibration-Free LIBS (CF-LIBS) eliminates the need for matrix-matched reference standards by deriving concentrations directly from plasma physics: the Boltzmann distribution governs excited-state populations, the Saha equation relates ionization-stage ratios to the electron temperature $T$ and electron density $n_e$, and a closure equation enforces the constraint that all elemental number fractions sum to unity \citep{Tognoni2010}.
Since its introduction by \citet{Ciucci1999}, CF-LIBS has been applied to geological surveying, industrial process control, cultural heritage analysis, and planetary science---most notably aboard the ChemCam and SuperCam instruments on NASA's Mars rovers \citep{Wiens2013,Maurice2021}.

Despite the elegance of the CF-LIBS approach, its computational demands are substantial.
A single inversion requires identifying emission lines against a database of $10^4$--$10^5$ atomic transitions, evaluating Voigt profiles for each identified line, performing weighted least-squares Boltzmann fits for every element, iterating a coupled Saha--Boltzmann system to self-consistency, and enforcing compositional closure.
When deployed in contexts that demand high throughput---online industrial monitoring, rapid survey campaigns, or manifold-based inference over grids of $10^6$--$10^8$ synthetic spectra---the CPU-bound nature of existing implementations becomes the primary bottleneck.
Real-time operation at the laser repetition rate (typically \SIrange{1}{20}{\hertz}, with $10$--$100$ accumulated spectra per analysis) requires processing hundreds of spectra per second, a regime where serial CPU code is inadequate.

The emergence of GPU computing in spectroscopy has demonstrated that massively parallel architectures can accelerate spectral calculations by orders of magnitude.
\citet{Kawahara2022} developed ExoJAX, a JAX-based framework for computing molecular and atomic cross-sections for exoplanet atmospheric retrieval, reporting throughput of approximately $10{,}000$ spectra per second on an NVIDIA V100 GPU.
\citet{Grimm2021} presented HELIOS-K, a CUDA~C++ opacity calculator for exoplanet atmospheres that achieves high throughput for Voigt profile evaluation across billions of spectral lines.
Both works demonstrate that GPU parallelism is well-suited to spectral computation, but neither addresses the specific algorithmic pipeline of CF-LIBS: multi-element Saha--Boltzmann coupling, Boltzmann plot fitting with variable line counts per element, compositional closure, and the iterative self-consistency loop that ties these components together.

To our knowledge, no GPU-accelerated implementation of the CF-LIBS pipeline has been reported.
This paper fills that gap by presenting a complete GPU-accelerated CF-LIBS library implemented in JAX \citep{Bradbury2018}, a Python framework for composable function transformations including automatic differentiation (\texttt{grad}), vectorization (\texttt{vmap}), and just-in-time compilation (\texttt{jit}) to GPU via the XLA compiler.
The choice of JAX over hand-tuned CUDA provides three advantages: (1)~automatic differentiation enables gradient-based joint optimization of temperature, electron density, and composition without manually deriving Jacobians; (2)~\texttt{vmap} provides batch parallelism with a single code path for both CPU and GPU; and (3)~Python-level development velocity is preserved while XLA delivers competitive GPU performance.

We contribute GPU-optimized implementations of five core algorithms:
\begin{enumerate}
  \item \textbf{Voigt profile evaluation} via the \citet{Weideman1994} $N=36$ rational approximation to the Faddeeva function, achieving a branch-free formulation compatible with JAX automatic differentiation (\Cref{sec:voigt}).
  \item \textbf{Vectorized Boltzmann fitting} using closed-form weighted least squares with a pad-and-mask strategy that batches variable-length line lists across elements without Python loops (\Cref{sec:boltzmann}).
  \item \textbf{Anderson-accelerated Saha--Boltzmann iteration} that reduces the iteration count of the electron density self-consistency loop by storing a history of previous iterates and solving a least-squares mixing problem \citep{Walker2011,Evans2018} (\Cref{sec:anderson}).
  \item \textbf{Softmax compositional closure} that reparameterizes the simplex constraint $\sum_i C_i = 1$, $C_i > 0$ as an unconstrained optimization in $\mathbb{R}^D$ with exact Jacobian, enabling gradient-based joint optimization (\Cref{sec:softmax}).
  \item \textbf{Batch forward model} that composes all four kernels into a vectorized pipeline mapping plasma parameters $(T, n_e, \bm{C})$ to synthetic spectra via \texttt{vmap}, enabling manifold generation at throughput exceeding $10{,}000$ spectra per second (\Cref{sec:batch}).
\end{enumerate}

All implementations use 64-bit floating-point arithmetic throughout to maintain the numerical precision required for spectroscopic analysis.
We benchmark each kernel individually and in an end-to-end pipeline on an NVIDIA Tesla V100S-PCIE-32GB GPU, demonstrating speedups ranging from $1.6\times$ (Anderson iteration reduction) to $76\times$ (Voigt profile throughput) with no loss of numerical accuracy relative to the CPU reference implementation.
We validate GPU--CPU numerical parity on 74~real LIBS spectra from the Aalto University Spectral Library and 6~synthetic ChemCam Calibration Target (CCCT) spectra.

The remainder of this paper is organized as follows.
\Cref{sec:methods} presents the mathematical formulation and GPU implementation strategy for each of the five kernels.
\Cref{sec:results} reports benchmark timings, accuracy validation, and real-data verification.
\Cref{sec:discussion} interprets the results in the context of prior GPU spectroscopy work and discusses limitations.
\Cref{sec:conclusion} summarizes the contributions and outlines future directions.
